\section{Temas adicionales: guía y calibración}
\subsection{Ingeniería de prompts de Self-Instruct}
Para que los prompts generados por Self-Instruct cubran a la vez matemáticas, código, razonamiento lógico y conversación, Weston recomienda escribir explícitamente en el prompt ``show your chain of thought'', ``highlight where you doubt'' y ``produce a verification statement''. Entre las prácticas concretas están:
\begin{itemize}
    \item Usar la plantilla: ``Step-by-step, explain ..., then verify with a judgement block.''.
    \item Agregar al final del prompt ``If you are unsure, state your uncertainty and ask for a cross-check''.
    \item Fijar de antemano el formato de la verified response (reasoning + final verdict), para que el entrenamiento del reward model pueda emparejarla con más facilidad.
\end{itemize}
\mbox{}

\subsection{Calibration y normalización de la recompensa}
Entrenar el meta judge requiere calibration, para poder juzgar con justicia entre la verified response y el draft. Weston nos recuerda que ``self evaluation is only going to help self-improve the model if inference compute for evaluation is controlled'' (SRT 3236-3252). Entre las estrategias habituales están reward clipping, temperature scaling, cross-check rate scheduling y el reentrenamiento del judge.

\begin{knowledgebox}{Calibration best practices}
Aplicar temperature scaling a los logits del judge, o elevar la frecuencia de cross-check cuando la verified response presenta baja confidence; aplicar clipping y normalization a los valores de reward, para evitar que un solo outlier haga que el self-rewarding se descontrole por completo.
\end{knowledgebox}

\subsection{Resumen de la sección}
El diseño de prompts y la calibration son condiciones necesarias para el éxito del self-rewarding: el primero garantiza la cobertura de las tareas sin perder el rumbo, y la segunda hace que el judge sea más estable durante el cross-check y no se desvíe por exceso de confianza ante una hallucination.
